\documentclass{zhvt-classic}

% 使用本地源流明体（GenRyuMin2 TW）按文件路径加载，不依赖系统字体名
% 版框（内框／界栏／书口侧栏／鱼尾）改用与版心印章同一印泥棕色：
%   印章图片 yshs.png 实色 = RGB(135,68,52) = HTML 874434
\usepackage{xcolor}
\usepackage{tikz}
% sealbrown 双注册：xcolor 供 TikZ/pgf，L3 供 cls 版框
\definecolor{sealbrown}{HTML}{874434}
\ExplSyntaxOn
\color_set:nnn { sealbrown } { HTML } { 874434 }
\ExplSyntaxOff

\zhvtset{
  grid color       = sealbrown,
  main font        = GenRyuMin2TW-SB,
  main font path   = C:/Users/Administrator/texmf/fonts/opentype/public/GenRyuMin2TW/,
jiazhu font      = GenRyuMin2TW-M,
  jiazhu font path = C:/Users/Administrator/texmf/fonts/opentype/public/GenRyuMin2TW/,
  % 页码框容纳 4 字：正常页「九十七」+前置空 1 格=4 行恰好；
  % 100+ 页「一百零一」共 5 行超出，末字「一/二/三」被框裁掉（2026-10-04 实测），故放宽到 5
  shukou page number chars = 5,
  % 书口书名框高从默认 6 字减到 4 字，把省出的 2 字高拨给 chapter 框，
  % 使 6 字书名+卷名放进 9 字高的 chapter 框后不再压线（2026-10-05）
  shukou title chars = 4,
}

\title{史记}[西南夷列传]
\author{汉·司马迁}
\maker{闲叟漫记·重製}


% ===== 索隐述赞样式（2026-11-08 四次定稿 · 愚大人亲手改定参数）=====
% ⚠ 竖排前提：自定义字体家族必须带 RawFeature={vertical:+vert}；
%   simhei.ttf 加 +vhal 会整体加载失败静默回退，只写 +vert；simfang 同理。
\newCJKfontfamily{\shuzanfont}[Path=C:/Windows/Fonts/,RawFeature={vertical:+vert}]{simhei.ttf}
\newCJKfontfamily{\shuzanbody}[Path=C:/Windows/Fonts/,RawFeature={vertical:+vert}]{simfang.ttf}
% 题头：152x30pt 圆角 6pt 棕底白字；baseline (0,2) 使整框落入界行格子不压版框线；
% 白字节点 x 不补偿、y +8.2pt（像素级墨迹实测 dx/dy≈0，L/R 留白约 20pt）。
\newcommand{\shuzanhead}[1]{%
  \begin{tikzpicture}[baseline={(0,2)}]
    \fill[sealbrown, rounded corners=6pt]
      (-76pt,-15pt) rectangle (76pt,15pt);
    \node[text=white, inner sep=0pt, anchor=center] at (0pt,8.2pt)
      {\shuzanfont\fontsize{23}{23}\selectfont #1};
  \end{tikzpicture}%
}
% ★ 修正书口卷名位置：chapter 框起始位置改回 cls 默认值，避免全局影响。
%   只在每个卷首页通过 \juan 命令里的空命令微调该页位置（安全，不会跨页污染）。
\ExplSyntaxOn
\AtBeginDocument{
  \dim_set:Nn \l__zhvt_shukou_chapter_title_start_dim
    { \l__zhvt_shukou_title_height_dim + \l__zhvt_font_size_dim * 2 }
  \dim_set:Nn \l__zhvt_shukou_page_number_start_dim
    { \l__zhvt_shukou_chapter_title_start_dim + \l__zhvt_shukou_chapter_title_height_dim }
}
\ExplSyntaxOff

\begin{document}

% ★★ 版心（书口）排布覆写：把书名放到鱼尾之下（2026-10-04）
%
%   cls 默认版心自上而下：书名（title 静态框，在鱼尾之上）→ 鱼尾 → 篇名（chapter 动态框）→ 页码
%   本模板改为：          鱼尾（位置不动）→ 书名 → 篇名 → 页码分割线 → 页码
%
%   手法：把 cls 放在 title 静态框里的书名，改放进 chapter 动态框（鱼尾之下）；
%         title 静态框留空。书名取 \l__zhvt_title_str，即 \title{...} 的第一参数，
%         所以换书时只改 \title，不必在此重复写书名。
%
%   ⚠ 三条必须遵守的写法（2026-10-04 实测）：
%     ① 汉字必须用 \hbox{...} 包住。flowfram 会把框内容装进一个宽 14.4mm（= 一个 line_sep
%        = 书口栏宽）的纵向 minipage 并按该宽度折行，而整页靠 PDF /Rotate 90 显示、
%        框内文字本身不变向 —— 不包 \hbox 时「围炉夜话」4 字会被逐字折成 4 行，
%        在版面上就成了横排 4 字并横向溢出版框。
%     ② 竖排下**一个 \hbox 就是一整列**。书名与篇名要同处一列，必须放进同一个 \hbox；
%        分成两个 \hbox 时，后一个会被当成新的一列、甩到隔壁栏去
%        （实测「序」跑到书口栏外侧 14.4mm 的正文列上）。中间的 \quad 在框内是横向空白，
%        显示出来正好是书名与篇名之间的纵向间隔。
%     ③ ★ 内容盒要写成 **与框等宽的定宽盒** `\hbox to \l__zhvt_shukou_chapter_title_height_dim{…\hfil}`。
%        实测：chapter 框是把内容盒沿纵向**居中**摆放的，内容宽度一变，书名就跟着漂：
%        自然宽度的 \hbox（书名 38.4mm）落在 x=157.2mm，加篇名后（57.6mm）落在 x=152.4mm，
%        目录页（无篇名）与正文页的书名就差半格。做成等宽盒后，无论有没有篇名，
%        内容固定从框顶起排，实测两页均落在 x=142.8mm，不再漂。
\ExplSyntaxOn
% 把书口 chapter 框（书名+篇名）起始位置下移 1 字高，避免 4 字卷名顶到鱼尾
% ⚠ cls 的 \__zhvt_calculate_dimensions: 在 \AtBeginDocument 里执行，会重算该 dim；
%   必须把覆写挂到 begin document 之后，否则会被二次覆盖（2026-10-05 实测）。
% (旧方案已移到导言区，直接修改 chapter 框起始位置)
% 卷标题命令 #1=书口篇名(鱼尾) #2=正文卷名（含卷序，卷名内空格用 ~）
%
% ★★ 目录页两项修订（2026-10-05，愚大人指出）★★
%   ① 页码起点统一对齐：
%      原状 —— 目录项写成 \par\hspace{4em}卷名\hspace{4em}页码，页码起点 = 卷名实宽 + 4em。
%      卷名实宽随字数变：「卷一~梁惠王上」6 字 +~、「卷七~离娄上」5 字 +~，
%      差 1 个字位 → 卷七~卷十 4 条的页码整体比其余 10 条高 27.21pt（一个字），看着上下不齐。
%      改法 —— 卷名补全角空格（宽 = 1 字 = 27.215pt）补足到 **7 字符位**再写进目录，
%      14 条页码起点遂全部落在同一条竖线上（起点对齐；末字随页码字数自然参差）。
%      ⚠ 只对目录里的副本 \l_tmpc_tl 补，正文卷名仍用 \l_tmpa_tl 原样 —— 正文像素不受影响。
%   ② 卷名整体上移半字位：前导 \hspace 由 4em 改 3.5em。
%      实测原状目录块在版框内偏下（天头留白 139.6pt / 地脚留白 110.7pt），
%      上移 0.5em = 13.6pt 后上下基本均分，视觉居中。
\NewDocumentCommand { \juan } { mm } {
	\clearpage
	\tl_set:Nn \l__zhvt_chapter_title_tl { #1 }
	\tl_set:Nn \l_tmpa_tl { #2 }
	\tl_set_eq:NN \l_tmpc_tl \l_tmpa_tl
	\int_step_inline:nn { 7 - \str_count:n { #2 } }
	  { \tl_put_right:Nn \l_tmpc_tl { 　 } }
	\tl_set:Nx \l_tmpb_tl { \zhnum { page } }
	\tl_put_left:Nn \l_tmpb_tl { \hspace { 4em } }
	\tl_put_left:Nn \l_tmpb_tl { \l_tmpc_tl }
	\tl_put_left:Nn \l_tmpb_tl { \par \hspace { 3.5em } }
	\addtocontents { toc } { \l_tmpb_tl }
	\skip_horizontal:n { \l__zhvt_font_size_dim * 4 }
	\l_tmpa_tl
	\par
	\thispagestyle { empty }
}
\cs_set_protected:Npn \__zhvt_set_shukou_content: {
	% ★ 鱼尾上方留白（title 静态框）嵌入闲叟钤印：
	%   整页 /Rotate 90，故此处把图用 \rotatebox{90} 转正为竖章；
	%   尺寸卡在版心条幅内（书内约 7mm 宽 × 33.8mm 高，留余量不压鱼尾）。
	\setstaticcontents * { title } {
		\vbox_to_ht:nn { \l__zhvt_baseline_skip_dim } {
			\vskip 14.2pt
			\hbox_to_wd:nn { \l__zhvt_shukou_title_height_dim } {
				\hfil
				\rotatebox [ origin = c ] { 90 } { \includegraphics [ width = 7mm , height = 33.8mm ] { yshs.png } }
				\hfil
			}
		}
	}
% ★ 版心（书口）统一改法（2026-11-16 愚大人定）：18pt＋框内居中＋书口文字同版框色
%   ① 原 27.2pt 时 7 字长篇名（如「廉颇蔺相如列传」）末字溢出 chapter 框、压住页码；
%      改 \fontsize{18pt}{18pt}，前后 \hfil 令字串在框内居中（不随篇名字数漂移）。
%   ② \vbox{\vskip 8.2pt} 定量下移，抵消缩字号后栏宽方向 8.2pt 上偏（像素实测 y 中心 543.3 = 原版）。
%      ⚠ \raisebox / \vspace* 在 flowfram 框内无效；换字号须同步换补偿：16pt→10.1、19pt→7.4。
%   ③ chapter 与 page~number 均置 \color{sealbrown}（HTML 874434），与版框／界行同色；正文仍黑。
%   ⚠ 页码字位、钤印、鱼尾、界行、正文一律未动；五部合订本同此口径。
	\setdynamiccontents * { chapter }
	{ \color{sealbrown} \vbox { \vskip 8.2pt \hbox to \l__zhvt_shukou_chapter_title_height_dim
		{ \hfil \fontsize{18pt}{18pt}\selectfont \l__zhvt_title_str \quad \l__zhvt_chapter_title_tl \hfil } } }
	\setdynamiccontents * { page~number } { \color{sealbrown} \quad \zhnum { page } }
}
\__zhvt_set_shukou_content:
\ExplSyntaxOff

\maketitle

% 单篇，无多卷目录，不渲染目录页

\juan{西南夷列传}{卷一百十六~西南夷列传第五十六}
\hspace*{2em}西南夷君长以什数。夜郎最大。其西靡莫之属以什数。滇最大。自滇以北君长以什数。邛都最大。此皆魋结。耕田。有邑聚。其外西自同师以东。北至楪榆。名为巂。昆明。皆编发。随畜迁徙。毋常处。毋君长。地方可数千里。自巂以东北。君长以什数。徙。筰都最大。自筰以东北。君长以什数。厓駹最大。其俗或士箸。或移徙。在蜀之西。自厓駹以东北。君长以什数。白马最大。皆氐类也。此皆巴蜀西南外蛮夷也。\\
\hspace*{2em}始楚威王时。使将军庄蹻将兵循江上。略巴。黔中以西。庄蹻者。故楚庄王苗裔也。蹻至滇池。方三百里。旁平地。肥饶数千里。以兵威定属楚。欲归报。会秦击夺楚巴。黔中郡。道塞不通。因还。以其众王滇。变服。从其俗。以长之。秦时常頞略通五尺道。诸此国颇置吏焉。十馀岁。秦灭。及汉兴。皆弃此国而开蜀故徼。巴蜀民或窃出商贾。取其筰马。僰僮。髦牛。以此巴蜀殷富。\\
\hspace*{2em}建元六年。大行王恢击东越。东越杀王郢以报。恢因兵威使番阳令唐蒙风指晓南越。南越食蒙蜀枸酱。蒙问所从来。曰﹁道西北牂柯。牂柯江广数里。出番禺城下﹂。蒙归至长安。问蜀贾人。贾人曰。﹁独蜀出枸酱。多持窃出市夜郎。夜郎者。临牂柯江。江广百馀步。足以行船。南越以财物役属夜郎。西至同师。然亦不能臣使也。﹂蒙乃上书说上曰。﹁南越王黄屋左纛。地东西万馀里。名为外臣。实一州主也。今以长沙。豫章往。水道多绝。难行。窃闻夜郎所有精兵。可得十馀万。浮船牂柯江。出其不意。此制越一奇也。诚以汉之彊。巴蜀之饶。通夜郎道。为置吏。易甚。﹂上许之。乃拜蒙为郎中将。将千人。食重万馀人。从巴蜀筰关入。遂见夜郎侯多同。蒙厚赐。喻以威德。约为置吏。使其子为令。夜郎旁小邑皆贪汉缯帛。以为汉道险。终不能有也。乃且听蒙约。还报。乃以为犍为郡。发巴蜀卒治道。自僰道指牂柯江。蜀人司马相如亦言西夷邛。筰可置郡。使相如以郎中将往喻。皆如南夷。为置一都尉。十馀县。属蜀。\\
\hspace*{2em}当是时。巴蜀四郡通西南夷道。戍转相饟。数岁。道不通。士罢饿离湿死者甚众。西南夷又数反。发兵兴击。秏费无功。上患之。使公孙弘往视问焉。还对。言其不便。及弘为御史大夫。是时方筑朔方以据河逐胡。弘因数言西南夷害。可且罢。专力事匈奴。上罢西夷。独置南夷夜郎两县一都尉。稍令犍为自葆就。\\
\hspace*{2em}及元狩元年。博望侯张骞使大夏来。言居大夏时见蜀布。邛竹。杖。使问所从来。曰﹁从东南身毒国。可数千里。得蜀贾人市﹂。或闻邛西可二千里有身毒国。骞因盛言大夏在汉西南。慕中国。患匈奴隔其道。诚通蜀。身毒国道便近。有利无害。於是天子乃令王然于。柏始昌。吕越人等。使间出西夷西。指求身毒国。至滇。滇王尝羌乃留。为求道西十馀辈。岁馀。皆闭昆明。莫能通身毒国。\\
\hspace*{2em}滇王与汉使者言曰。﹁汉孰与我大。﹂及夜郎侯亦然。以道不通故。各自以为一州主。不知汉广大。使者还。因盛言滇大国。足事亲附。天子注意焉。\\
\hspace*{2em}及至南越反。上使驰义侯因犍为发南夷兵。且兰君恐远行。旁国虏其老弱。乃与其众反。杀使者及犍为太守。汉乃发巴蜀罪人尝击南越者八校尉击破之。会越已破。汉八校尉不下。即引兵还。行诛头兰。头兰。常隔滇道者也。已平头兰。遂平南夷为牂柯郡。夜郎侯始倚南越。南越已灭。会还诛反者。夜郎遂入朝。上以为夜郎王。\\
\hspace*{2em}南越破后。及汉诛且兰。邛君。并杀筰侯。厓駹皆振恐。请臣置吏。乃以邛都为越巂郡。筰都为沈犁郡。厓駹为汶山郡。广汉西白马为武都郡。\\
\hspace*{2em}上使王然于以越破及诛南夷兵威风喻滇王入朝。滇王者。其众数万人。其旁东北有劳浸。靡莫。皆同姓相扶。未肯听。劳浸。靡莫数侵犯使者吏卒。元封二年。天子发巴蜀兵击灭劳浸。靡莫。以兵临滇。滇王始首善。以故弗诛。滇王离难西南夷。举国降。请置吏入朝。於是以为益州郡。赐滇王王印。复长其民。\\
\hspace*{2em}西南夷君长以百数。独夜郎。滇受王印。滇小邑。最宠焉。\\
\hspace*{2em}太史公曰。楚之先岂有天禄哉。在周为文王师。封楚。及周之衰。地称五千里。秦灭诸侯。唯楚苗裔尚有滇王。汉诛西南夷。国多灭矣。唯滇复为宠王。然南夷之端。见枸酱番禺。大夏杖。邛竹。西夷后揃。剽分二方。卒为七郡。\\

% ===== 索隐述赞（司马贞《史记索隐》）· 统一样式 =====
% 题头「索隐述赞曰」五字并入反白框；正文每四句一列，定宽盒 \hfil 居列中；
% \fontsize{24}{41} 列距对齐全书界行（121+41k）；\clearpage 使述赞独占一页。
\clearpage
\shuzanhead{索隐述赞曰}\linebreak
{\shuzanbody\fontsize{24}{41}\selectfont
\hbox to \hsize{\hfil 西南外徼。庄蹻首通。汉因大夏。乃命唐蒙。\hfil}\linebreak
\hbox to \hsize{\hfil 劳洸。靡莫。异俗殊风。夜郎最大。\hfil}\linebreak
\hbox to \hsize{\hfil 邛。筰称雄。及置郡县。万代推功。\hfil}%
\par}
\end{document}
